\chapter{2009年广东卷（文）}

\section{单选题}

\begin{problem}
已知全集 \(U = \R\)，则正确表示集合 \(M = \{-1, 0, 1\}\) 和 \(N = \{x \mid x^2 + x = 0\}\) 关系的韦恩（Venn）图是（\quad）

\bitmapfigure[max width=0.88\linewidth]{img/2009/guangdong_liberal/q1_options.png}
\end{problem}

\begin{answer}
B
\end{answer}

\begin{solution}
由 \(x^2+x=x(x+1)=0\)，得 \(N=\{-1,0\}\). 因为 \(M=\{-1,0,1\}\)，所以 \(N\subset M\)，且 \(1\in M\) 而 \(1\notin N\). 因此 \(N\) 应位于 \(M\) 的内部，故选 B.
\end{solution}

\begin{problem}
下列 \(n\) 的取值中，使 \(\i^n = 1\)（\(\i\) 是虚数单位）的是（\quad）

\choices
    {\(n = 2\)}
    {\(n = 3\)}
    {\(n = 4\)}
    {\(n = 5\)}
\end{problem}

\begin{answer}
C
\end{answer}

\begin{solution}
虚数单位的幂以 \(4\) 为周期，且 \(\i^2=-1\)，\(\i^3=-\i\)，\(\i^4=1\)，\(\i^5=\i\). 给出的四个取值中只有 \(n=4\) 时 \(\i^n=1\)，故选 C.
\end{solution}

\begin{problem}
已知平面向量 \( \bs{a} = (x, 1) \)，\( \bs{b} = (-x, x^{2}) \)，则向量 \( \bs{a} + \bs{b} \)（\quad）

\choices
    {平行于 \(x\) 轴}
    {平行于第一、三象限的角平分线}
    {平行于 \(y\) 轴}
    {平行于第二、四象限的角平分线}
\end{problem}

\begin{answer}
C
\end{answer}

\begin{solution}
\(\bs{a}+\bs{b}=(x-x,1+x^2)=(0,1+x^2)\). 由于 \(1+x^2>0\)，该向量非零且横坐标为 \(0\)，所以它平行于 \(y\) 轴，故选 C.
\end{solution}

\begin{problem}
若函数 \( y = f(x) \) 是函数 \( y = a^x \)（\(a > 0\)，且 \(a \neq 1\)）的反函数，且 \(f(2) = 1\)，则 \(f(x) =\)（\quad）

\choices
    {\(\log_2 x\)}
    {\(\frac{1}{2^x}\)}
    {\(\log_{\frac{1}{2}} x\)}
    {\(2^{x - 2}\)}
\end{problem}

\begin{answer}
A
\end{answer}

\begin{solution}
函数 \(y=a^x\) 的反函数为 \(f(x)=\log_a x\). 由 \(f(2)=1\)，得 \(\log_a2=1\)，从而 \(a=2\). 因此 \(f(x)=\log_2x\)，故选 A.
\end{solution}

\begin{problem}
已知等比数列 \(\{a_n\}\) 的公比为正数，且 \(a_3 \cdot a_9 = 2a_5^2\)，\(a_2 = 1\)，则 \(a_1 =\)（\quad）

\choices
    {\(\frac{1}{2}\)}
    {\(\frac{\sqrt{2}}{2}\)}
    {\(\sqrt{2}\)}
    {\(2\)}
\end{problem}

\begin{answer}
B
\end{answer}

\begin{solution}
设等比数列的首项为 \(a_1\)，公比为 \(q\). 由题意，\(q>0\)，且 \(a_2=a_1q=1\). 又
\(a_3a_9=(a_1q^2)(a_1q^8)=a_1^2q^{10}\)，\(2a_5^2=2(a_1q^4)^2=2a_1^2q^8\).

因为 \(a_1\ne0\)，所以 \(q^2=2\). 结合 \(q>0\)，得 \(q=\sqrt2\)，从而 \(a_1=1/q=\frac{\sqrt2}{2}\)，故选 B.
\end{solution}

\begin{problem}
给定下列四个命题：
\begin{circlelist}
\item 若一个平面内的两条直线与另一个平面都平行，那么这两个平面相互平行；
\item 若一个平面经过另一个平面的垂线，那么这两个平面相互垂直；
\item 垂直于同一直线的两条直线相互平行；
\item 若两个平面垂直，那么一个平面内与它们的交线不垂直的直线与另一个平面也不垂直.
\end{circlelist}
其中，为真命题的是（\quad）

\choices
    {\circled{1}和\circled{2}}
    {\circled{2}和\circled{3}}
    {\circled{3}和\circled{4}}
    {\circled{2}和\circled{4}}
\end{problem}

\begin{answer}
D
\end{answer}

\begin{solution}
命题\(\circled{1}\)不成立. 例如取平面 \(\alpha\colon z=0\) 和平面 \(\beta\colon y=0\)，在 \(\alpha\) 内取直线 \(l_1\colon y=1\)，\(z=0\) 和 \(l_2\colon y=2\)，\(z=0\). 两条直线都不与 \(\beta\) 相交，因而都平行于 \(\beta\)，但 \(\alpha\) 与 \(\beta\) 相交，不能推出两平面平行. 该命题缺少“两条直线相交”的条件.

命题\(\circled{2}\)成立. 若平面 \(\alpha\) 经过直线 \(l\)，且 \(l\perp\beta\)，那么平面 \(\alpha\) 中有一条直线垂直于平面 \(\beta\). 由两平面垂直的判定定理，得 \(\alpha\perp\beta\).

命题\(\circled{3}\)不成立. 例如取 \(z\) 轴为公共直线，取 \(x\) 轴和 \(y\) 轴为另外两条直线；两轴都垂直于 \(z\) 轴，但 \(x\) 轴与 \(y\) 轴相交而不平行.

命题\(\circled{4}\)成立. 设两平面的交线为 \(l\)，其中一个平面内的直线为 \(m\). 两平面垂直时，该平面内垂直于另一个平面的直线必垂直于交线 \(l\)；因此若 \(m\) 不垂直于 \(l\)，则 \(m\) 不可能垂直于另一个平面. 所以真命题为 \(\circled{2}\) 和 \(\circled{4}\)，故选 D.
\end{solution}

\begin{problem}
已知 \(\triangle ABC\) 中，\(\angle A\)，\(\angle B\)，\(\angle C\) 的对边分别为 \(a\)，\(b\)，\(c\). 若 \(a = c = \sqrt{6} + \sqrt{2}\)，且 \(\angle A = 75^{\circ}\)，则 \(b =\)（\quad）

\choices
    {\(2\)}
    {\(4 + 2\sqrt{3}\)}
    {\(4 - 2\sqrt{3}\)}
    {\(\sqrt{6} -\sqrt{2}\)}
\end{problem}

\begin{answer}
A
\end{answer}

\begin{solution}
因为 \(a=c\)，所以 \(\angle A=\angle C=75^\circ\)，从而 \(\angle B=180^\circ-75^\circ-75^\circ=30^\circ\). 由正弦定理，

\(\frac{b}{\sin30^\circ}=\frac{a}{\sin75^\circ}\).

又 \(a=\sqrt6+\sqrt2=4\sin75^\circ\)，所以 \(b=4\sin75^\circ\cdot\frac{\sin30^\circ}{\sin75^\circ}=2\)，故选 A.
\end{solution}

\begin{problem}
函数 \( f(x) = (x - 3)\e^x \) 的单调递增区间是（\quad）

\choices
    {\((- \infty, 2)\)}
    {\((0,3)\)}
    {\((1,4)\)}
    {\((2, +\infty)\)}
\end{problem}

\begin{answer}
D
\end{answer}

\begin{solution}
\(f(x)=(x-3)\e^x\)，所以

\(f'(x)=\e^x+(x-3)\e^x=(x-2)\e^x\).

由于 \(\e^x>0\)，当 \(x>2\) 时 \(f'(x)>0\)，函数单调递增，故选 \((2,+\infty)\)，即 D.
\end{solution}

\begin{problem}
函数 \( y = 2\cos^2 \left( x - \frac{\pi}{4} \right) - 1 \) 是（\quad）

\choices
    {最小正周期为 \(\pi\) 的奇函数}
    {最小正周期为 \(\pi\) 的偶函数}
    {最小正周期为 \(\frac{\pi}{2}\) 的奇函数}
    {最小正周期为 \(\frac{\pi}{2}\) 的偶函数}
\end{problem}

\begin{answer}
A
\end{answer}

\begin{solution}
利用 \(2\cos^2u-1=\cos2u\)，得

\(y=\cos\left(2x-\frac{\pi}{2}\right)=\sin2x\).

函数 \(\sin2x\) 是奇函数，且最小正周期为 \(\pi\)，故选 A.
\end{solution}

\begin{problem}
广州 2010 年亚运会火炬传递在 \(A\)，\(B\)，\(C\)，\(D\)，\(E\) 五个城市之间进行，各城市之间的距离（单位：百公里）见下表. 若以 \(A\) 为起点，\(E\) 为终点，每个城市经过且只经过一次，那么火炬传递的最短路线距离是（\quad）

\begin{center}
\begin{tabular}{|c|c|c|c|c|c|}
\hline
 & \(A\) & \(B\) & \(C\) & \(D\) & \(E\) \\
\hline
\(A\) & 0 & 5 & 4 & 5 & 6 \\
\hline
\(B\) & 5 & 0 & 7 & 6 & 2 \\
\hline
\(C\) & 4 & 7 & 0 & 9 & 8.6 \\
\hline
\(D\) & 5 & 6 & 9 & 0 & 5 \\
\hline
\(E\) & 6 & 2 & 8.6 & 5 & 0 \\
\hline
\end{tabular}
\end{center}

\choices
    {\(20.6\)}
    {\(21\)}
    {\(22\)}
    {\(23\)}
\end{problem}

\begin{answer}
B
\end{answer}

\begin{solution}
从 \(A\) 到 \(E\) 需依次经过 \(B\)，\(C\)，\(D\) 的某一种排列. 六条路线的长度分别为

\(A\to B\to C\to D\to E\)：\(5+7+9+5=26\)，
\(A\to B\to D\to C\to E\)：\(5+6+9+8.6=28.6\)，
\(A\to C\to B\to D\to E\)：\(4+7+6+5=22\)，
\(A\to C\to D\to B\to E\)：\(4+9+6+2=21\)，
\(A\to D\to B\to C\to E\)：\(5+6+7+8.6=26.6\)，
\(A\to D\to C\to B\to E\)：\(5+9+7+2=23\).

最短长度为 \(21\)，故选 B.
\end{solution}

\section{填空题}

\begin{problem}
某篮球队 6 名主力队员在最近三场比赛中投进的三分球个数如下表所示：
\begin{center}
\begin{tabular}{|c|c|c|c|c|c|c|}
\hline
队员 \(i\) & 1 & 2 & 3 & 4 & 5 & 6 \\
\hline
三分球个数 & \(a_1\) & \(a_2\) & \(a_3\) & \(a_4\) & \(a_5\) & \(a_6\) \\
\hline
\end{tabular}
\end{center}
如图是统计该 6 名队员在最近三场比赛中投进的三分球总数的程序框图，则图中判断框应填\fillinblank{}，输出的 \(s=\)\fillinblank{}.（注：框图中的赋值符号“=”也可以写成“←”或“：=”）

\begingroup
\problemresetlayout
\bitmapfigure[width=0.32\linewidth]{img/2009/guangdong_liberal/q11_fig1.png}
\endgroup
\end{problem}

\begin{answer}
\(i\leqslant 6\)，

\(a_1+a_2+a_3+a_4+a_5+a_6\)
\end{answer}

\begin{solution}
程序初始化为 \(s=0\)，\(i=1\). 第一次到达判断框时尚未执行 \(i=i+1\)，若判断为“是”，就把 \(a_i\) 加入 \(s\)，再执行 \(i=i+1\) 返回判断框. 为了依次累加 \(a_1\)，\(a_2\)，\(\ldots\)，\(a_6\)，判断框在 \(i=1\)，\(2\)，\(\ldots\)，\(6\) 时都应为“是”，在 \(i=7\) 时应为“否”，故判断条件为 \(i\leqslant6\).

循环中的累加过程为
\(s=0+a_1\)，\(s=a_1+a_2\)，\(\ldots\)，\(s=a_1+a_2+a_3+a_4+a_5+a_6\). 因而输出的 \(s=a_1+a_2+a_3+a_4+a_5+a_6\)，两个空依次为 \(i\leqslant6\) 和 \(a_1+a_2+a_3+a_4+a_5+a_6\).
\end{solution}

\begin{problem}
某单位 200 名职工的年龄分布情况如图，现要从中抽取 40 名职工作样本.
用系统抽样法，将全体职工随机按 \(1 \sim 200\) 编号，并按编号顺序平均分为 40 组（\(1 \sim 5\) 号，\(6 \sim 10\) 号，\(\cdots\)，\(196 \sim 200\) 号）. 若第 5 组抽出的号码为 22，则第 8 组抽出的号码应是\fillinblank{}. 若用分层抽样方法，则 \(40\text{ 岁}\)以下年龄段应抽取\fillinblank{}人.
\bitmapfigure[max width=0.34\linewidth]{img/2009/guangdong_liberal/q12_fig1.png}
\end{problem}

\begin{answer}
37，20
\end{answer}

\begin{solution}
每组有 \(200\div40=5\) 个编号. 第 \(5\) 组的编号为 \(21\sim25\)，其中抽出的号码为 \(22\)，即抽取的是该组的第 \(2\) 个号码. 系统抽样中各组抽取相同位置的号码，因此第 \(8\) 组抽出的号码为
\(22+(8-5)\times5=37\).

由年龄分布图知，\(40\text{ 岁}\)以下职工占 \(50\%\). 按分层抽样方法，抽取的 \(40\) 人中该年龄段应有
\(40\times50\%=20\) 人.
\end{solution}

\begin{problem}
以点 \((2,-1)\) 为圆心且与直线 \(x + y = 6\) 相切的圆的方程是\fillinblank{}.
\end{problem}

\begin{answer}
\((x-2)^2+(y+1)^2=\frac{25}{2}\)
\end{answer}

\begin{solution}
圆心为 \(P(2,-1)\)，圆心到直线 \(x+y-6=0\) 的距离为
\(r=\frac{|2+(-1)-6|}{\sqrt{1^2+1^2}}=\frac{5}{\sqrt2}\).
由于圆与该直线相切，半径就是这个距离，所以 \(r^2=\frac{25}{2}\). 因此圆的方程为
\((x-2)^2+(y+1)^2=\frac{25}{2}\).
\end{solution}

\begin{problem}
若直线 \(\begin{cases}
x = 1 - 2t,\\
y = 2 + 3t,
\end{cases}\)（\(t\) 为参数）与直线 \(4x + ky = 1\) 垂直，则常数 \(k =\)\fillinblank{}.

\end{problem}

\begin{answer}
\(-6\)
\end{answer}

\begin{solution}
参数方程中的直线方向向量为 \(\bs{v}_1=(-2,3)\). 直线 \(4x+ky=1\) 的一个方向向量为 \(\bs{v}_2=(k,-4)\). 两直线垂直，所以
\(\bs{v}_1\cdot\bs{v}_2=(-2)k+3(-4)=-2k-12=0\).
解得 \(k=-6\).
\end{solution}

\begin{problem}
如图，点 \(A\)，\(B\)，\(C\) 是圆 \(O\) 上的点，且 \(AB = 4\)，\(\angle ACB = 30^{\circ}\)，则圆 \(O\) 的面积等于\fillinblank{}.

\FigureLayoutDeclare{img/2009/guangdong_liberal/q15_fig1.png}{4.5cm}{36bp 39bp 264bp 11bp}
\bitmapfigure[max width=0.36\linewidth]{img/2009/guangdong_liberal/q15_fig1.png}
\end{problem}

\begin{answer}
\(16\pi\)
\end{answer}

\begin{solution}
圆周角 \(\angle ACB=30^\circ\) 所对的圆心角为 \(60^\circ\). 设圆 \(O\) 的半径为 \(R\)，则弦长公式给出
\(AB=2R\sin30^\circ\).
代入 \(AB=4\)，得 \(4=2R\cdot\frac12\)，所以 \(R=4\). 因此圆 \(O\) 的面积为 \(\pi R^2=16\pi\).
\end{solution}

\section{解答题}

\begin{problem}
已知向量 \( \bs{a} = (\sin\theta, -2) \) 与 \( \bs{b} = (1, \cos\theta) \) 互相垂直，其中 \( \theta \in \left(0, \frac{\pi}{2}\right) \).
\begin{enumerate}
\item 求 \(\sin\theta\cos\theta\) 的值；
\item 若 \(5\cos(\theta - \varphi) = 3\sqrt{5}\cos\varphi\)，\(0 < \varphi < \frac{\pi}{2}\)，求 \(\cos\varphi\) 的值.
\end{enumerate}
\end{problem}

\begin{solution}
\begin{enumerate}
\item 因为 \(\bs{a}\) 与 \(\bs{b}\) 互相垂直，所以
\(\bs{a}\cdot\bs{b}=\sin\theta-2\cos\theta=0\)，即 \(\sin\theta=2\cos\theta\). 又因为 \(\theta\in\left(0,\frac{\pi}{2}\right)\)，所以 \(\sin\theta>0\)，\(\cos\theta>0\)，且
\(\sin^2\theta+\cos^2\theta=1\). 代入 \(\sin\theta=2\cos\theta\)，得 \(5\cos^2\theta=1\). 因此
\(\cos\theta=\frac{1}{\sqrt5}=\frac{\sqrt5}{5}\)，\(\sin\theta=\frac{2}{\sqrt5}=\frac{2\sqrt5}{5}\)，从而 \(\sin\theta\cos\theta=\frac25\).
\item 由 \(5\cos(\theta-\varphi)=3\sqrt5\cos\varphi\) 及余弦差公式，得
\(5(\cos\theta\cos\varphi+\sin\theta\sin\varphi)=3\sqrt5\cos\varphi\). 代入上一问的结果，得
\(\sqrt5\cos\varphi+2\sqrt5\sin\varphi=3\sqrt5\cos\varphi\)，从而 \(\sin\varphi=\cos\varphi\). 由于 \(0<\varphi<\frac{\pi}{2}\)，所以 \(\varphi=\frac{\pi}{4}\)，故
\(\cos\varphi=\frac{\sqrt2}{2}\).
\end{enumerate}
\end{solution}

\begin{problem}
某高速公路收费站入口处的安全标识墩如图 1 所示. 墩的上半部分是正四棱锥 \(P-EFGH\)，下半部分是长方体 \(ABCD-EFGH\). 图 2，图 3 分别是该标识墩的正（主）视图和俯视图（单位：\(\mathrm{~cm}\)）.

\begin{enumerate}
\item 请画出该安全标识墩的侧（左）视图；
\item 求该安全标识墩的体积；
\item 证明：直线 \(BD \perp\) 平面 \(PEG\).
\end{enumerate}

\bitmapfigure[max width=0.72\linewidth]{img/2009/guangdong_liberal/q17_fig1.png}
\end{problem}

\begin{solution}
\begin{enumerate}
\item 由正视图和俯视图可知，底部是底边长 \(40\mathrm{~cm}\)、高 \(20\mathrm{~cm}\) 的矩形，上部正四棱锥的底边长为 \(40\mathrm{~cm}\)，高为 \(60\mathrm{~cm}\). 正四棱锥的顶点在正方形底面的中心正上方，所以从左侧投影时，顶点仍落在底边中点的正上方. 因此侧视图是一个底边长 \(40\mathrm{~cm}\)、高 \(20\mathrm{~cm}\) 的矩形，其上叠加一个底边长 \(40\mathrm{~cm}\)、高 \(60\mathrm{~cm}\) 的等腰三角形.
\item 长方体的体积为
\(V_1=40\times40\times20=32000\mathrm{~cm}^3\). 正四棱锥的体积为
\(V_2=\frac13\times40\times40\times60=32000\mathrm{~cm}^3\). 因此标识墩的体积为
\(V=V_1+V_2=64000\mathrm{~cm}^3\).
\item 设 \(O\) 为正方形 \(EFGH\) 的中心，则 \(O\) 是对角线 \(EG\) 与 \(FH\) 的交点. 因为 \(P-EFGH\) 是正四棱锥，故 \(PO\perp\) 平面 \(EFGH\)，从而 \(PO\perp FH\). 又正方形的两条对角线互相垂直，所以 \(EG\perp FH\). 直线 \(PO\) 与 \(EG\) 是平面 \(PEG\) 内相交的两条直线，且 \(FH\) 同时垂直于它们，因此 \(FH\perp\) 平面 \(PEG\). 上、下两个正方形平行且对应边方向相同，故 \(BD\parallel FH\)，于是 \(BD\perp\) 平面 \(PEG\).
\end{enumerate}
\end{solution}

\begin{problem}
随机抽取某中学甲，乙两班各 10 名同学，测量他们的身高（单位：\(\mathrm{~cm}\)），获得身高数据的茎叶图如图.

\begin{enumerate}
\item 根据茎叶图判断哪个班的平均身高较高；
\item 计算甲班的样本方差；
\item 现从乙班这 10 名同学中随机抽取两名身高不低于 \(173 \mathrm{~cm}\) 的同学，求身高为 \(176 \mathrm{~cm}\) 的同学被抽中的概率.
\end{enumerate}

\begin{center}
\begin{tabular}{r|c|l}
甲班 & & 乙班 \\
\hline
\(2\) & 18 & \(1\) \\
\(9\)，\(9\)，\(1\)，\(0\) & 17 & \(0\)，\(3\)，\(6\)，\(8\)，\(9\) \\
\(8\)，\(8\)，\(3\)，\(2\) & 16 & \(2\)，\(5\)，\(8\) \\
\(8\) & 15 & \(9\)
\end{tabular}
\end{center}
\end{problem}

\begin{solution}
\begin{enumerate}
\item 由茎叶图，甲班的身高数据为 \(182\)，\(179\)，\(179\)，\(171\)，\(170\)，\(168\)，\(168\)，\(163\)，\(162\)，\(158\)，其平均数为
\(\bar{x}_{\text{甲}}=\frac{1700}{10}=170\). 乙班的身高数据为 \(181\)，\(179\)，\(178\)，\(176\)，\(173\)，\(170\)，\(168\)，\(165\)，\(162\)，\(159\)，其平均数为
\(\bar{x}_{\text{乙}}=\frac{1711}{10}=171.1\). 因为 \(171.1>170\)，所以乙班的平均身高较高.
\item 甲班的样本方差为
\[
s^2=\frac1{10}\bigl[12^2+9^2+9^2+1^2+0^2+2^2+2^2+7^2+8^2+12^2\bigr]
=\frac{572}{10}=57.2
\]
\item 乙班身高不低于 \(173 \mathrm{~cm}\) 的同学有 \(173\)，\(176\)，\(178\)，\(179\)，\(181\)，共 \(5\) 人. 任取其中 \(2\) 人，所有等可能的取法共有 \(\mathrm C_5^2\) 种；含有身高 \(176 \mathrm{~cm}\) 同学的取法共有 \(\mathrm C_4^1\) 种. 因此所求概率为
\(\frac{\mathrm C_4^1}{\mathrm C_5^2}=\frac{4}{10}=\frac25\).
\end{enumerate}
\end{solution}

\begin{problem}
已知椭圆 \(G\) 的中心在坐标原点，长轴在 \(x\) 轴上，离心率为 \(\frac{\sqrt{3}}{2}\)，两个焦点分别为 \(F_{1}\) 和 \(F_{2}\)，椭圆 \(G\) 上一点到 \(F_{1}\) 和 \(F_{2}\) 的距离之和为 12. 圆 \(C_{k}: x^{2} + y^{2} + 2kx - 4y - 21 = 0\)（\(k \in \R\)）的圆心为点 \(A_{k}\).
\begin{enumerate}
\item 求椭圆 \(G\) 的方程；
\item 求 \( \triangle A_{k}F_{1}F_{2} \) 面积；
\item 问是否存在圆 \(C_{k}\) 包围椭圆 \(G\)？请说明理由.
\end{enumerate}
\end{problem}

\begin{solution}
\begin{enumerate}
\item 设椭圆 \(G\) 的长半轴、短半轴和半焦距分别为 \(a\)，\(b\)，\(c\). 由椭圆定义，椭圆上一点到两个焦点的距离之和为 \(2a\)，所以 \(2a=12\)，得 \(a=6\). 又由离心率 \(e=\frac ca=\frac{\sqrt3}{2}\)，得 \(c=3\sqrt3\). 在椭圆中 \(b^2=a^2-c^2\)，所以 \(b^2=36-27=9\). 由于长轴在 \(x\) 轴上，椭圆方程为
\(\displaystyle\frac{x^2}{36}+\frac{y^2}{9}=1\).
\item 将圆 \(C_k\) 的方程配方，得 \((x+k)^2+(y-2)^2=k^2+25\)，
所以圆心 \(A_k=(-k,2)\)，焦点 \(F_1\)，\(F_2\) 在 \(x\) 轴上，且 \(F_1F_2=2c=6\sqrt3\). 三角形 \(A_kF_1F_2\) 的底边为 \(F_1F_2\)，高为点 \(A_k\) 到 \(x\) 轴的距离 \(2\)，故
\(\left[ \triangle A_kF_1F_2\right]=\frac12\times6\sqrt3\times2=6\sqrt3\).
\item 若圆 \(C_k\) 包围椭圆 \(G\)，则椭圆的两个顶点 \((6,0)\)，\((-6,0)\) 都在圆内或圆上. 由圆的方程，点 \((6,0)\) 到圆心 \(A_k\) 的距离平方为 \((k+6)^2+4\)，于是
\((k+6)^2+4\leq k^2+25\)，即 \(k\leq-\frac54\). 同理，点 \((-6,0)\) 到圆心的距离平方为 \((k-6)^2+4\)，从而
\((k-6)^2+4\leq k^2+25\)，即 \(k\geq\frac54\). 两个条件不能同时成立，所以不存在这样的圆 \(C_k\).
\end{enumerate}
\end{solution}

\begin{problem}
已知点 \(\left(1, \frac{1}{3}\right)\) 是函数 \(f(x) = a^x\)（\(a > 0\)，\(a \neq 1\)）的图象上一点. 等比数列 \(\{a_n\}\) 的前 \(n\) 项和为 \(f(n) - c\). 数列 \(\{b_n\}\)（\(b_n > 0\)）的首项为 \(c\)，且前 \(n\) 项和 \(S_n\) 满足 \(S_n - S_{n-1} = \sqrt{S_n} + \sqrt{S_{n-1}}\)（\(n \ge 2\)）.

\begin{enumerate}
\item 求数列 \(\{a_n\}\) 和 \(\{b_n\}\) 的通项公式；
\item 若数列 \(\left\{\frac{1}{b_n b_{n+1}}\right\}\) 的前 \(n\) 项和为 \(T_n\)，问满足 \(T_n > \frac{1000}{2009}\) 的最小正整数 \(n\) 是多少？
\end{enumerate}
\end{problem}

\begin{solution}
\begin{enumerate}
\item 因为点 \(\left(1,\frac13\right)\) 在 \(y=a^x\) 的图象上，所以 \(a=\frac13\)，即 \(f(n)=\left(\frac13\right)^n\). 对 \(n\geq2\)，由前 \(n\) 项和相减，得
\(a_n=(f(n)-c)-(f(n-1)-c)=\left(\frac13\right)^n-\left(\frac13\right)^{n-1}=-\frac{2}{3^n}\). 因此 \(a_2=-\frac29\). 另一方面 \(a_1=f(1)-c=\frac13-c\)，而等比数列公比为 \(\frac13\)，所以 \(a_2=\frac13a_1\)，从而 \(a_1=-\frac23\)，解得 \(c=1\).

因此 \(a_n=-\frac{2}{3^n}\)（\(n\geqslant1\)）. 又 \(b_1=c=1\)，所以 \(S_1=1\). 对 \(n\geq2\)，由 \(b_n=S_n-S_{n-1}\) 及题设 \(S_n-S_{n-1}=\sqrt{S_n}+\sqrt{S_{n-1}}\)，得
令 \(u_n=\sqrt{S_n}\)，则 \(u_n>0\)，且
\((u_n-u_{n-1})(u_n+u_{n-1})=u_n+u_{n-1}\). 因为 \(u_n+u_{n-1}>0\)，所以 \(u_n-u_{n-1}=1\). 由 \(u_1=1\)，归纳得 \(u_n=n\)，即 \(S_n=n^2\). 因此
\(b_n=S_n-S_{n-1}=n^2-(n-1)^2=2n-1\).
\item 由上一问，\(b_n=2n-1\)，所以
\[
T_n=\sum_{k=1}^n\frac1{(2k-1)(2k+1)}
=\frac12\sum_{k=1}^n\left(\frac1{2k-1}-\frac1{2k+1}\right)
=\frac12\left(1-\frac1{2n+1}\right)=\frac{n}{2n+1}
\]
题设不等式等价于
\(\frac{n}{2n+1}>\frac{1000}{2009}\). 由于分母均为正，交叉相乘得 \(2009n>2000n+1000\)，即 \(9n>1000\)，所以 \(n>\frac{1000}{9}\). 满足条件的最小正整数为 \(n=112\).
\end{enumerate}
\end{solution}

\begin{problem}
已知二次函数 \(y = g(x)\) 的导函数的图象与直线 \(y = 2x\) 平行，且 \(y = g(x)\) 在 \(x = -1\) 处取得极小值 \(m - 1\)（\(m \neq 0\)）. 设 \(f(x) = \frac{g(x)}{x}\).

\begin{enumerate}
\item 若曲线 \(y = f(x)\) 上的点 \(P\) 到点 \(Q(0,2)\) 的距离的最小值为 \(\sqrt{2}\)，求 \(m\) 的值；
\item \(k\)（\(k \in \R\)）如何取值时，函数 \(y = f(x) - kx\) 存在零点，并求出零点.
\end{enumerate}
\end{problem}

\begin{solution}
\begin{enumerate}
\item 因为 \(g'(x)\) 的图象与 \(y=2x\) 平行，所以可设 \(g'(x)=2x+b\). 又 \(g(x)\) 在 \(x=-1\) 处取得极小值，故 \(g'(-1)=0\)，从而 \(b=2\). 积分得 \(g(x)=x^2+2x+c\). 由 \(g(-1)=m-1\)，得 \(c=m\)，所以
\(g(x)=x^2+2x+m\)，\(f(x)=\frac{g(x)}x=x+2+\frac{m}{x}\)，其中 \(x\ne0\).

点 \(P(x,f(x))\) 到 \(Q(0,2)\) 的距离平方为
\[
d^2=x^2+\left(f(x)-2\right)^2=x^2+\left(x+\frac{m}{x}\right)^2
=2x^2+2m+\frac{m^2}{x^2}
\]
令 \(u=x^2>0\)，则
\(d^2=2u+2m+\frac{m^2}{u}\). 由基本不等式，
\(2u+\frac{m^2}{u}\geq2\sqrt{2m^2}=2\sqrt2|m|\)，当 \(u=\frac{|m|}{\sqrt2}\) 时等号成立，所以
\(d_{\min}^2=2m+2\sqrt2|m|\). 题设 \(d_{\min}=\sqrt2\)，故
\(m+\sqrt2|m|=1\).
当 \(m>0\) 时，\(m(1+\sqrt2)=1\)，得 \(m=\sqrt2-1\)；当 \(m<0\) 时，\(m(1-\sqrt2)=1\)，得 \(m=-\sqrt2-1\).
\item 由 \(f(x)-kx=0\) 且 \(x\ne0\)，得 \((1-k)x^2+2x+m=0\).
当 \(k=1\) 时，方程化为 \(2x+m=0\)，由于 \(m\ne0\)，唯一零点为 \(x=-\frac m2\). 当 \(k\ne1\) 时，判别式为
\(\Delta=4-4m(1-k)=4[1-m(1-k)]\). 因而有零点的充要条件是
\(1-m(1-k)\geqslant0\)，零点为
\(\displaystyle x=\frac{1\pm\sqrt{1-m(1-k)}}{k-1}\). 因为 \(m\ne0\)，这些根不可能为 \(0\)，均满足 \(x\ne0\)；等号成立时两个根重合为 \(x=-m\).

结合第一问的结果，当 \(m=\sqrt2-1>0\) 时，条件化为 \(k\geqslant1-\frac1m=-\sqrt2\)；当 \(m=-\sqrt2-1<0\) 时，条件化为 \(k\leqslant1-\frac1m=\sqrt2\). 因此，\(m=\sqrt2-1\) 时取 \(k\geqslant-\sqrt2\)，\(m=-\sqrt2-1\) 时取 \(k\leqslant\sqrt2\)；不满足对应条件时不存在零点.
\end{enumerate}
\end{solution}
